\documentclass[12pt,a4paper]{exam}
\usepackage[utf8]{inputenc}
\usepackage{amsmath,amssymb,amsfonts}
\usepackage{geometry}
\usepackage{enumitem}
\usepackage{graphicx}
\usepackage{hyperref}
\usepackage{xcolor}
\geometry{a4paper, top=2cm, bottom=2cm, left=2.5cm, right=2.5cm}
\hypersetup{colorlinks=true, linkcolor=blue, urlcolor=blue}
\renewcommand{\thequestion}{\textbf{\arabic{question}}}
\renewcommand{\questionlabel}{\thequestion.}
\pointname{ marks}
\pointsdroppedatright
\bracketedpoints
\setlength{\parindent}{0pt}
\setlength{\emergencystretch}{3em}
\allowdisplaybreaks
\newsavebox{\schmathbox}
\newcommand{\fitmath}[1]{%
  \sbox{\schmathbox}{$\displaystyle #1$}%
  \ifdim\wd\schmathbox>\linewidth
    \resizebox{\linewidth}{!}{\usebox{\schmathbox}}%
  \else
    \usebox{\schmathbox}%
  \fi}

\begin{document}

\thispagestyle{empty}
\begin{center}
  \textbf{\LARGE Diag} \\[0.3cm]
  \textbf{Time Allowed: 3 Hours} \hfill \textbf{Maximum Marks: 80} \\[0.3cm]
  \rule{\textwidth}{0.5pt}
\end{center}

\vspace{0.3cm}

\noindent\textbf{General Instructions:}
\begin{enumerate}
  \item {\normalfont\normalcolor This question paper contains 40 questions. All questions are compulsory.}
  \item {\normalfont\normalcolor Questions are grouped into sections A through E.}
  \item {\normalfont\normalcolor Use of calculators is NOT permitted.}
\end{enumerate}
\bigskip
\noindent\textbf{\large Section A: Multiple Choice \& Assertion-Reason (1 mark each)}

\begin{questions}
\question[1] A $1.5$ m tall boy is standing at some distance from a $30$ m tall building. The angle of elevation from his eyes to the top of the building is $45^\circ$. The height above his eye level is:
\begin{enumerate}[label=(\Alph*)]
  \item (A) $30$ m
  \item (B) $28.5$ m
  \item (C) $31.5$ m
  \item (D) $1.5$ m
\end{enumerate}
\vspace{0.15cm}

\question[1] The shadow of a tower standing on a level ground is found to be $40$ m longer when the Sun's altitude is $30^\circ$ than when it is $60^\circ$. The height of the tower is not required, but if we consider a simpler case where the shadow is $h\sqrt{3}$ when angle is $30^\circ$, what is the height $h$ if the shadow is $30$ m?
\begin{enumerate}[label=(\Alph*)]
  \item (A) $30\sqrt{3}$ m
  \item (B) $30$ m
  \item (C) $10\sqrt{3}$ m
  \item (D) $10$ m
\end{enumerate}
\vspace{0.15cm}

\question[1] If the angle between two tangents drawn from a point $P$ to a circle of radius $r$ and center $O$ is $60^\circ$, then the length of $OP$ is:
\begin{enumerate}[label=(\Alph*)]
  \item (A) r
  \item (B) 2r
  \item $(C) r\sqrt{3}$
  \item (D) r/2
\end{enumerate}
\vspace{0.15cm}

\question[1] A tangent $PQ$ at a point $P$ of a circle of radius $5 \text{ cm}$ meets a line through the center $O$ at a point $Q$ so that $OQ = 13 \text{ cm}$. The length of $PQ$ is:
\begin{enumerate}[label=(\Alph*)]
  \item (A) 10 cm
  \item (B) 8 cm
  \item (C) 12 cm
  \item (D) 7 cm
\end{enumerate}
\vspace{0.15cm}

\question[1] The number of tangents that can be drawn from a point lying inside a circle is:
\begin{enumerate}[label=(\Alph*)]
  \item (A) 0
  \item (B) 1
  \item (C) 2
  \item (D) Infinite
\end{enumerate}
\vspace{0.15cm}

\question[1] If two tangents inclined at an angle of $80^\circ$ are drawn to a circle of radius $3 \text{ cm}$, then the angle between the two radii at the center is:
\begin{enumerate}[label=(\Alph*)]
  \item $(A) 80^\circ$
  \item $(B) 90^\circ$
  \item $(C) 100^\circ$
  \item $(D) 120^\circ$
\end{enumerate}
\vspace{0.15cm}

\question[1] The length of a tangent from a point $A$ at a distance of $10 \text{ cm}$ from the center of the circle is $8 \text{ cm}$. The radius of the circle is:
\begin{enumerate}[label=(\Alph*)]
  \item (A) 4 cm
  \item (B) 6 cm
  \item (C) 7 cm
  \item (D) 9 cm
\end{enumerate}
\vspace{0.15cm}

\question[1] Assertion (A): The angle of elevation of the top of a tower from a point on the ground, which is $30 \text{ m}$ away from the foot of the tower, is $45^\circ$. The height of the tower is $30 \text{ m}$. Reason (R): In a right-angled triangle, $\tan(\theta) = \frac{\text{Opposite}}{\text{Adjacent}}$.
\begin{enumerate}[label=(\Alph*)]
  \item (A) Both Assertion (A) and Reason (R) are true and Reason (R) is the correct explanation of Assertion (A)
  \item (B) Both Assertion (A) and Reason (R) are true but Reason (R) is NOT the correct explanation of Assertion (A)
  \item (C) Assertion (A) is true but Reason (R) is false
  \item (D) Assertion (A) is false but Reason (R) is true
\end{enumerate}
\vspace{0.15cm}

\question[1] Assertion (A): If the length of the shadow of a vertical pole is equal to its height, then the angle of elevation of the Sun is $45^\circ$. Reason (R): The angle of elevation is the angle formed by the line of sight with the horizontal when the object is above the horizontal level.
\begin{enumerate}[label=(\Alph*)]
  \item (A) Both Assertion (A) and Reason (R) are true and Reason (R) is the correct explanation of Assertion (A)
  \item (B) Both Assertion (A) and Reason (R) are true but Reason (R) is NOT the correct explanation of Assertion (A)
  \item (C) Assertion (A) is true but Reason (R) is false
  \item (D) Assertion (A) is false but Reason (R) is true
\end{enumerate}
\vspace{0.15cm}

\question[1] Assertion (A): The length of a kite string making an angle of $30^\circ$ with the ground to reach a height of $50 \text{ m}$ is $100 \text{ m}$. Reason (R): $\sin(\theta) = \frac{\text{Opposite}}{\text{Hypotenuse}}$.
\begin{enumerate}[label=(\Alph*)]
  \item (A) Both Assertion (A) and Reason (R) are true and Reason (R) is the correct explanation of Assertion (A)
  \item (B) Both Assertion (A) and Reason (R) are true but Reason (R) is NOT the correct explanation of Assertion (A)
  \item (C) Assertion (A) is true but Reason (R) is false
  \item (D) Assertion (A) is false but Reason (R) is true
\end{enumerate}
\vspace{0.15cm}

\question[1] Assertion (A): The angle of depression from a point $10 \text{ m}$ above sea level to a ship at sea is $30^\circ$. The distance of the ship from the base of the point is $10\sqrt{3} \text{ m}$. Reason (R): Angle of depression is equal to the angle of elevation by alternate interior angles.
\begin{enumerate}[label=(\Alph*)]
  \item (A) Both Assertion (A) and Reason (R) are true and Reason (R) is the correct explanation of Assertion (A)
  \item (B) Both Assertion (A) and Reason (R) are true but Reason (R) is NOT the correct explanation of Assertion (A)
  \item (C) Assertion (A) is true but Reason (R) is false
  \item (D) Assertion (A) is false but Reason (R) is true
\end{enumerate}
\vspace{0.15cm}

\question[1] Assertion (A): If a ladder $10 \text{ m}$ long is placed against a wall such that it makes an angle of $60^\circ$ with the ground, then the foot of the ladder is $5 \text{ m}$ away from the wall. Reason (R): $\cos(\theta) = \frac{\text{Adjacent}}{\text{Hypotenuse}}$.
\begin{enumerate}[label=(\Alph*)]
  \item (A) Both Assertion (A) and Reason (R) are true and Reason (R) is the correct explanation of Assertion (A)
  \item (B) Both Assertion (A) and Reason (R) are true but Reason (R) is NOT the correct explanation of Assertion (A)
  \item (C) Assertion (A) is true but Reason (R) is false
  \item (D) Assertion (A) is false but Reason (R) is true
\end{enumerate}
\vspace{0.15cm}

\question[1] Assertion (A): The length of a tangent drawn from an external point to a circle is equal to the length of another tangent drawn from the same point. Reason (R): The lengths of tangents drawn from an external point to a circle are equal.
\begin{enumerate}[label=(\Alph*)]
  \item (A) Both Assertion (A) and Reason (R) are true and Reason (R) is the correct explanation of Assertion (A)
  \item (B) Both Assertion (A) and Reason (R) are true but Reason (R) is NOT the correct explanation of Assertion (A)
  \item (C) Assertion (A) is true but Reason (R) is false
  \item (D) Assertion (A) is false but Reason (R) is true
\end{enumerate}
\vspace{0.15cm}

\question[1] Assertion (A): The tangent at any point of a circle is perpendicular to the radius through the point of contact. Reason (R): The perpendicular distance from the center of the circle to the tangent is equal to the radius.
\begin{enumerate}[label=(\Alph*)]
  \item (A) Both Assertion (A) and Reason (R) are true and Reason (R) is the correct explanation of Assertion (A)
  \item (B) Both Assertion (A) and Reason (R) are true but Reason (R) is NOT the correct explanation of Assertion (A)
  \item (C) Assertion (A) is true but Reason (R) is false
  \item (D) Assertion (A) is false but Reason (R) is true
\end{enumerate}
\vspace{0.15cm}

\question[1] Assertion (A): The number of tangents that can be drawn from a point inside a circle is zero. Reason (R): A point inside a circle lies in the interior, and any line passing through it will be a secant, not a tangent.
\begin{enumerate}[label=(\Alph*)]
  \item (A) Both Assertion (A) and Reason (R) are true and Reason (R) is the correct explanation of Assertion (A)
  \item (B) Both Assertion (A) and Reason (R) are true but Reason (R) is NOT the correct explanation of Assertion (A)
  \item (C) Assertion (A) is true but Reason (R) is false
  \item (D) Assertion (A) is false but Reason (R) is true
\end{enumerate}
\vspace{0.15cm}

\question[1] Assertion (A): If the angle between two tangents drawn from a point P to a circle of radius $r$ and center O is $60^\circ$, then $OP = 2r$. Reason (R): The line segment from the center to the external point bisects the angle between the tangents.
\begin{enumerate}[label=(\Alph*)]
  \item (A) Both Assertion (A) and Reason (R) are true and Reason (R) is the correct explanation of Assertion (A)
  \item (B) Both Assertion (A) and Reason (R) are true but Reason (R) is NOT the correct explanation of Assertion (A)
  \item (C) Assertion (A) is true but Reason (R) is false
  \item (D) Assertion (A) is false but Reason (R) is true
\end{enumerate}
\vspace{0.15cm}

\question[1] Assertion (A): The length of the tangent from a point $A$ at a distance of $5 \text{ cm}$ from the center of the circle is $4 \text{ cm}$. The radius of the circle is $3 \text{ cm}$. Reason (R): In a right triangle formed by radius, tangent, and hypotenuse (distance to center), $r^2 + t^2 = d^2$.
\begin{enumerate}[label=(\Alph*)]
  \item (A) Both Assertion (A) and Reason (R) are true and Reason (R) is the correct explanation of Assertion (A)
  \item (B) Both Assertion (A) and Reason (R) are true but Reason (R) is NOT the correct explanation of Assertion (A)
  \item (C) Assertion (A) is true but Reason (R) is false
  \item (D) Assertion (A) is false but Reason (R) is true
\end{enumerate}
\vspace{0.15cm}

\end{questions}
\bigskip
\noindent\textbf{\large Section B: Very Short Answer (2 marks each)}

\begin{questions}
\question[2] A ladder $15 \text{ m}$ long makes an angle of $60^\circ$ with the wall. Find the distance of the foot of the ladder from the wall.
\vspace{0.15cm}

\question[2] The angle of elevation of the top of a tower from a point on the ground, which is $30 \text{ m}$ away from the foot of the tower, is $30^\circ$. Find the height of the tower.
\vspace{0.15cm}

\question[2] A kite is flying at a height of $60 \text{ m}$ above the ground. The string attached to the kite is temporarily tied to a point on the ground. The inclination of the string with the ground is $60^\circ$. Find the length of the string.
\vspace{0.15cm}

\question[2] From the top of a $7 \text{ m}$ high building, the angle of elevation of the top of a cable tower is $60^\circ$. If the distance between the building and the tower is $7\sqrt{3} \text{ m}$, find the height of the tower.
\vspace{0.15cm}

\question[2] A $1.5 \text{ m}$ tall boy is standing at some distance from a $30 \text{ m}$ tall building. The angle of elevation from his eyes to the top of the building increases from $30^\circ$ to $60^\circ$ as he walks towards the building. Find the distance he walked.
\vspace{0.15cm}

\question[2] A tangent $PQ$ at a point $P$ of a circle of radius $5$ cm meets a line through the center $O$ at a point $Q$ so that $OQ = 13$ cm. Find the length of $PQ$.
\vspace{0.15cm}

\question[2] If tangents $PA$ and $PB$ from a point $P$ to a circle with center $O$ are inclined to each other at an angle of $80^\circ$, then find $\angle POA$.
\vspace{0.15cm}

\question[2] Two concentric circles are of radii $5$ cm and $3$ cm. Find the length of the chord of the larger circle which touches the smaller circle.
\vspace{0.15cm}

\question[2] A quadrilateral $\text{ABCD}$ is drawn to circumscribe a circle. If $AB = 6$ cm, $BC = 7$ cm, and $CD = 4$ cm, find the length of $AD$.
\vspace{0.15cm}

\question[2] If a point $P$ is $26$ cm away from the center $O$ of a circle and the length of the tangent $PT$ drawn from $P$ to the circle is $10$ cm, find the radius of the circle.
\vspace{0.15cm}

\end{questions}
\bigskip
\noindent\textbf{\large Section C: Short Answer (3 marks each)}

\begin{questions}
\question[3] A vertical tower stands on a horizontal plane and is surmounted by a vertical flagstaff of height $6$ m. At a point on the plane, the angle of elevation of the bottom and the top of the flagstaff are $30^{\circ}$ and $60^{\circ}$ respectively. Find the height of the tower.
\vspace{0.15cm}

\question[3] A $1.5$ m tall boy is standing at some distance from a $30$ m tall building. The angle of elevation from his eyes to the top of the building increases from $30^{\circ}$ to $60^{\circ}$ as he walks towards the building. Find the distance he walked towards the building.
\vspace{0.15cm}

\question[3] The angle of elevation of the top of a lighthouse from a boat is $30^{\circ}$. If the boat moves $100$ m towards the lighthouse, the angle of elevation becomes $45^{\circ}$. Find the height of the lighthouse.
\vspace{0.15cm}

\question[3] A kite is flying at a height of $60$ m above the ground. The string attached to the kite is temporarily tied to a point on the ground. The inclination of the string with the ground is $60^{\circ}$. Find the length of the string, assuming that there is no slack in the string.
\vspace{0.15cm}

\question[3] From a point on the ground, the angles of elevation of the bottom and the top of a transmission tower fixed at the top of a $20$ m high building are $45^{\circ}$ and $60^{\circ}$ respectively. Find the height of the tower.
\vspace{0.15cm}

\question[3] A quadrilateral $\text{ABCD}$ is drawn to circumscribe a circle. Prove that $AB + CD = AD + BC$.
\vspace{0.15cm}

\question[3] From a point $Q$, the length of the tangent to a circle is $24$ cm and the distance of $Q$ from the centre is $25$ cm. Find the radius of the circle.
\vspace{0.15cm}

\question[3] Two concentric circles are of radii $5$ cm and $3$ cm. Find the length of the chord of the larger circle which touches the smaller circle.
\vspace{0.15cm}

\question[3] If tangents $PA$ and $PB$ from a point $P$ to a circle with centre $O$ are inclined to each other at an angle of $80^\circ$, then find $\angle POA$.
\vspace{0.15cm}

\end{questions}
\bigskip
\noindent\textbf{\large Section D: Long Answer (4-5 marks each)}

\begin{questions}
\question[5] A vertical tower stands on a horizontal plane and is surmounted by a vertical flagstaff of height $h$. At a point on the plane, the angles of elevation of the bottom and the top of the flagstaff are $\alpha$ and $\beta$ respectively. Prove that the height of the tower is $\frac{h \tan \alpha}{\tan \beta - \tan \alpha}$.
\vspace{0.15cm}

\question[5] From a point on the ground, the angles of elevation of the bottom and the top of a transmission tower fixed at the top of a $20 \text{ m}$ high building are $45^\circ$ and $60^\circ$ respectively. Find the height of the tower.
\vspace{0.15cm}

\question[5] The angle of elevation of a cloud from a point $60 \text{ m}$ above a lake is $30^\circ$ and the angle of depression of the reflection of the cloud in the lake is $60^\circ$. Find the height of the cloud from the surface of the lake.
\vspace{0.15cm}

\question[5] Two poles of equal heights are standing opposite each other side by side and opposite each other on either side of the road, which is $80 \text{ m}$ wide. From a point between them on the road, the angles of elevation of the top of the poles are $60^\circ$ and $30^\circ$ respectively. Find the height of the poles and the distances of the point from the poles.
\vspace{0.15cm}

\end{questions}
\bigskip

\end{document}
